\begin{figure}[tb]
\centering
\begin{tikzpicture}[font=\footnotesize,>=Latex,line cap=round,
  bank/.style={very thick},
  barrier/.style={densely dashed,gray!80},
  fluid/.style={fill=blue!8}]
  \begin{scope}[xshift=0.3cm]
    \fill[blue!8] (0,0) circle (2.65);
    \fill[black!9]
      (2,0) .. controls (1.25,.25) and (.25,1.25) .. (0,2)
      .. controls (-.25,1.25) and (-1.25,.25) .. (-2,0)
      .. controls (-1.25,-.25) and (-.25,-1.25) .. (0,-2)
      .. controls (.25,-1.25) and (1.25,-.25) .. cycle;
    \draw[barrier] (0,0) circle (2);
    \draw[bank]
      (2,0) .. controls (1.25,.25) and (.25,1.25) .. (0,2);
    \draw[bank,gray!55]
      (0,2) .. controls (-.25,1.25) and (-1.25,.25) .. (-2,0);
    \draw[bank]
      (-2,0) .. controls (-1.25,-.25) and (-.25,-1.25) .. (0,-2);
    \draw[bank,gray!55]
      (0,-2) .. controls (.25,-1.25) and (1.25,-.25) .. (2,0);
    \foreach \p in {(2,0),(0,2),(-2,0),(0,-2)}
      \fill \p circle (2pt);
    \draw[gray!75] (0,0) circle (.58);
    \node[align=center,font=\scriptsize] at (0,0) {starting\\side};
    \node at (1.15,1.15) {$L_1$};
    \node at (-1.15,1.15) {$L_2$};
    \node at (-1.15,-1.15) {$L_3$};
    \node at (1.15,-1.15) {$L_4$};
    \node[align=center] at (0,2.97) {unqueried annular back};
    \draw[->,thin] (.78,2.87)--(1.15,2.18);
    \node[align=left,anchor=west] at (2.8,.95)
      {four-tip\\frontier};
    \draw[->,thin] (2.78,.77)--(1.43,.30);
    \node[align=left,anchor=west] at (2.8,-.80)
      {open attachment\\interval};
    \draw[->,thin] (2.78,-.70)--(1.74,-.97);
    \node[align=center] at (0,-3.05)
      {Dashed circle: stopping barrier.\\
       Its open intervals lie inside the glued domain.};
  \end{scope}

  \begin{scope}[xshift=8.85cm,yshift=.20cm]
    \node[align=center] at (0,2.28) {Two distinct banks at a contact};
    \fill[blue!8]
      (-1.55,1.70)--(1.55,1.70)--(1.55,.85)
      .. controls (.85,.85) and (.50,0) .. (0,0)
      .. controls (-.50,0) and (-.85,.85) .. (-1.55,.85)--cycle;
    \fill[blue!8]
      (-1.55,-1.70)--(1.55,-1.70)--(1.55,-.85)
      .. controls (.85,-.85) and (.50,0) .. (0,0)
      .. controls (-.50,0) and (-.85,-.85) .. (-1.55,-.85)--cycle;
    \draw[bank]
      (-1.55,.85) .. controls (-.85,.85) and (-.50,0) .. (0,0)
      .. controls (.50,0) and (.85,.85) .. (1.55,.85);
    \draw[bank,gray!55]
      (-1.55,-.85) .. controls (-.85,-.85) and (-.50,0) .. (0,0)
      .. controls (.50,0) and (.85,-.85) .. (1.55,-.85);
    \fill (0,0) circle (2.2pt);
    \node at (0,1.30) {first access};
    \node at (0,-1.30) {second access};
    \draw[->,thin] (0,.96)--(0,.20);
    \draw[->,thin] (0,-.96)--(0,-.20);
    \node[align=center] at (0,-2.30)
      {The contact is boundary,\\not a gluing interval.};
  \end{scope}
\end{tikzpicture}
\caption{The decomposition at an inner four-tip frontier. The four
monochromatic frontier arcs, drawn alternately black and gray, bound
Jordan lobes $L_1,\ldots,L_4$ with the stopping barrier. Each lobe is
attached to the untouched annular back along an open barrier interval;
the common shading denotes the resulting unqueried domain. The starting
side remains in the complement. At the other end the construction is
reversed. The inset records two accesses at a possible contact between
different banks: the contact is not added to the fluid and does not join
the accesses. All curves are schematic. Actual banks may be rough and
may touch each other, although each individual frontier arc is simple.}
\label{fig:frontier-lobes}
\end{figure}
